\section*{Resumen}

El presente trabajo de investigación abordó la problemática de la ciberseguridad en infraestructuras académicas frente a amenazas cibernéticas autopropagables. El objetivo principal consistió en desarrollar una simulación basada en agentes empleando la plataforma GAMA para analizar la dinámica de propagación de un ataque de ransomware inspirado en WannaCry en una red local universitaria y evaluar una estrategia de contención. La metodología empleada incluyó el modelado espacial de siete laboratorios de computación interconectados mediante archivos vectoriales de QGIS, y la implementación de reglas probabilísticas de transmisión lateral basadas en la exposición de puertos y niveles de actualización del sistema operativo. Los resultados demostraron una rápida saturación de la infraestructura en ausencia de controles. No obstante, tras activar un protocolo de contingencia automatizado basado en el aislamiento lógico al alcanzar un umbral crítico del treinta por ciento de equipos infectados, se consiguió proteger el setenta por ciento restante de la red. Finalmente, se concluyó que la topología de la red influyó de forma directa en la velocidad de dispersión de la amenaza y que la segmentación y el aislamiento temprano resultaron altamente eficientes como defensas activas.